\section{Introduction}
\label{sec:intro}

% \carlos{Right now, the main proposition is that SWE-bench training data is kind of hard to acquire $\rightarrow$ SWE-smith makes it easy to make training data for agents. We should mention more about why we want to do this. I think we should mention the fact that improvements of proprietary chat models is probably the \textit{most} significant source of improvements (e.g. claude 3.5 was a major improvement, even on existing scaffolds), however, of course this is limited by the fact they aren't open source and thus inaccessible for researchers to experiment with more deeply. SWE-smith is our proposed solution, and we show that this strategy makes it easier to train LMs for SWE-agent. Once results are available, I think we should generally adapt the current intro section to emphasize the significance of LM training for SWE scaffolds. SWE-smith makes it possible to train agent models!}

Language models (LMs) have demonstrated exciting promise towards automating software engineering (SE) tasks, as indicated by continued improvement in performance on benchmarks for real world GitHub issue resolution~\citep{jimenez_swe-bench_2024,yang_swe-bench_2024}, code generation with libraries~\citep{zhao_commit0_2024,zhuo2024bigcodebenchbenchmarkingcodegeneration}, test generation~\citep{mündler2025swtbenchtestingvalidatingrealworld}, and execution environment construction~\citep{bogin2024superevaluatingagentssetting,eliseeva2025envbenchbenchmarkautomatedenvironment,vergopoulos2025automatedbenchmarkgenerationrepositorylevel}.
SWE-bench~\citep{jimenez_swe-bench_2024} has emerged as a core evaluation for measuring improvements in LMs for SE.
Much of this progress thus far has come from better inference frameworks for SE tasks, such as agents~\citep{yang_swe-agent_2024,wang_openhands_2024,antoniades_swe-search_2024,su2025learnbyinteractdatacentricframeworkselfadaptive} and pipeline based methods~\citep{xia_agentless_2024,orwall2024moatless,zhang_autocoderover_2024}.
However, efforts to directly improve base LMs on SE tasks remain limited, largely due to the lack of high-quality training data and software engineering environments.

Recent contributions have attempted to construct training datasets for SWE-bench by either crawling large quantities of pull requests (PRs)~\citep{jimenez_swe-bench_2024,wei2025swerladvancingllmreasoning,xie2025swefixertrainingopensourcellms} or re-applying SWE-bench style collection to a distinct set of GitHub repositories~\citep{badertdinov2024scaling,pan_training_2024} \kilian{aren't SWE-bench style collections an example for crawling PRs?-- so less of an \emph{either or} }.
\kilian{I'd rather continue with 'however, in order to generate training trajectories, execution environments have to be set up for every single instance, \dots }PR-based datasets do not necessarily provide execution environments, rendering them unserviceable for work towards interactive, agent based approaches at scale.
SWE-bench style training sets have several apparent drawbacks.
First, a key step of collecting SWE-bench task instances is determining installation and testing specifications for multiple versions of a codebase across time. \kilian{The fact that we have to set it up for multiple versions of the codebase is very important, so should make sure reader understands that we need to do this potentially for every instance, because every instance is ankered to a different commit -- because this is why swe-smith is so much more efficient: install once, generate unlimited}
This may entail intensive human labor to manually write these specs and debug issues such as rediscovering undocumented installation steps or toggling versioning for dependencies~\citep{yang_swe-bench_2024} \kilian{what does toggling versioning mean?}
~\cite{pan_training_2024} reports that collecting $2,438$ instances from $11$ repositories required $200$ human annotation hours.
Each task instance's environment also needs significant storage, at about $2.6$GB per instance, totaling $6$TB~\citep{pan_training_2024}.
Collectively, these challenges make scaling up SWE-bench style collection challenging due to labor and storage costs.

\kilian{Last paragraph presented the problem; I think here we should directly lead with the new idea "set up environment once for one revision, then \emph{generate} an arbitrary amount of \emph{synthetic issues} from it (and only then add technical details) -- I think having an immediate 'ah, that's how they want to solve the problem' moment is more satisfying to the reader}
To address this problem, we introduce \bugs{}, a SWE-bench style training dataset \carlos{(We should maybe say something different from ``SWE-bench style'' since we just said ``SWE-bench style'' training datasets were bad in the 3rd sentence of above paragraph.)} of $50$k task instances with execution environments, sourced from $128$ real-world Python repositories. \diyi{To some extent, i think the contribution can be framed as a methodology to generate scalable training dataset? framing it as a dataset seems a undersell to me.}
To achieve such scale, we \textit{invert} SWE-bench's collection strategy of crawling task instances first and creating environments second.
Inspired by automated environment setup work~\citep{bogin2024superevaluatingagentssetting}, we first define environments for a large number of repositories.
Then, we employ $4$ classes of methods to create $100$s to $1000$s of software bugs per repository.
Using \bugs{}, we can scale up the training data for SE agents (e.g., task instances, repositories, environments) at a fraction of prior human labor and storage costs.

Our experiments with \bugs{} demonstrate the benefits of scaling up agentic data for SE.
We create a dataset of $10$k trajectories generated by running SWE-agent~\citep{yang_swe-agent_2024} with Claude 3.7 Sonnet~\citep{anthropicclaude37} on \bugs{} task instances.
\john{TO DO: Results}
